Il presente capitolo descrive l'implementazione in MATLAB del simulatore
presentato nel capitolo precedente. L'attenzione è rivolta al passaggio dai
vincoli progettuali alle strutture dati, alle funzioni e agli aggiornamenti
di stato che producono il programma di esecuzione delle missioni.
La descrizione fa riferimento al percorso effettivamente attivato da
\texttt{main.m}: assegnazione offline degli ordini, suddivisione in batch
statici, esecuzione con \texttt{colored\allowbreak Petri} e valutazione delle deadline
mediante una previsione euclidea congelata al primo avvio riuscito.

La persistenza riguarda il lavoro cumulato, le posizioni delle risorse,
gli storici e le deadline. La rete e la marcatura vengono invece ricostruite
per ogni batch. I moduli del motore temporale alternativo presenti nel
progetto non sono richiamati dal flusso principale; di conseguenza, finestre
di osservazione, tempi residui e controllo a orizzonte mobile non descrivono
il funzionamento corrente.

\section{Organizzazione del software}
\label{sec:impl-organizzazione}

Lo script principale individua la propria directory e aggiunge al percorso
di ricerca MATLAB le cartelle dei moduli utilizzati. I nomi del file Excel
e delle immagini sono risolti rispetto a tale directory, evitando di far
dipendere la lettura degli ingressi dalla cartella di lavoro corrente.
La Tabella~\ref{tab:impl-moduli} riassume la ripartizione delle responsabilità.

\begin{table}[htbp]
	\centering
	\small
	\begin{tabularx}{\linewidth}{@{}p{4.1cm}X@{}}
		\toprule
		Modulo & Responsabilità nel flusso attivo \\
		\midrule
		\texttt{main.m} & Configurazione, coordinamento dei batch e riepiloghi. \\
		\texttt{import\allowbreak Excel} & Lettura delle colonne e normalizzazione dell'ordine. \\
		\texttt{model\allowbreak Generation} & Compatibilità, proprietari degli ordini, matrici e archi. \\
		\texttt{batching} & Suddivisione della lista mantenendo gli ordini interi. \\
		\texttt{colored\_\allowbreak petri\_\allowbreak nets3} & Scatto delle transizioni, marcatura e cronologia. \\
		\texttt{optimization} & Risoluzione dei conflitti e rami storici di blocco delle risorse. \\
		\texttt{path\allowbreak Plan} & Mappa, coordinate, raggiungibilità e A*. \\
		\texttt{utility} & Preparazione, aggiornamenti, deadline e grafici. \\
		\bottomrule
	\end{tabularx}
	\caption{Organizzazione dei moduli richiamati dalla simulazione.}
	\label{tab:impl-moduli}
\end{table}

Il nome della cartella \texttt{optimization} non implica che il programma
risolva un problema di ottimizzazione globale. Con il proprietario già fissato,
la risoluzione dei conflitti deve rispettare il colore assegnato all'ordine.
Analogamente, \texttt{plot\allowbreak Temporal\allowbreak Results} è una funzione grafica usata anche
nel flusso statico: il nome non implica l'attivazione del motore a finestre.

Le cartelle \texttt{tests}, \texttt{docs} e \texttt{tmp} contengono
rispettivamente verifiche, documentazione e materiali di lavoro. La cartella
\texttt{temporal} conserva il motore alternativo. Questi componenti vanno
distinti dalle dipendenze necessarie per eseguire lo script principale.

\subsection{Configurazione dell'esperimento}

La configurazione imposta otto risorse, la soglia di blocco a infinito e
la visualizzazione dei grafici attiva. Il riferimento di otto ore e la durata
nominale di 180 secondi vengono passati alla divisione in batch, ma la
possibilità di ridurre automaticamente il numero di risorse è disabilitata.
Pertanto non introducono un arresto del motore al termine del turno.

Per le deadline sono configurati un fattore di percorrenza di 5.5 celle al
secondo, una maggiorazione del 20\%, un contributo accessorio di 60 secondi
e un margine aggiuntivo nullo per missione. Il campo
\texttt{Distance\allowbreak Metric} riporta il valore euclideo; la funzione invocata
calcola direttamente norme euclidee. La modifica della sola stringa non
seleziona un diverso algoritmo di stima.

Gli ingressi della configurazione corrente sono
\texttt{excel\_\allowbreak files/1su11.xlsx} e le immagini
\texttt{map1\_\allowbreak 1\allowbreak Filled.png} e \texttt{map1\_\allowbreak 1.png} nella cartella
\texttt{path\allowbreak Plan}. L'esecuzione richiede che MATLAB disponga delle funzioni
per la lettura del foglio e delle classi \texttt{binary\allowbreak Occupancy\allowbreak Map} e
\texttt{planner\allowbreak AStar\allowbreak Grid}. La versione dell'ambiente e la disponibilità di
queste dipendenze fanno parte delle condizioni di riproducibilità.

\subsubsection{Estratto di implementazione}

Il Listing~\ref{lst:impl-configurazione} riporta la configurazione effettivamente
utilizzata dallo script principale. In particolare, i parametri della previsione
temporale sono separati sia dal limite di otto ore sia dai parametri del batching.

\begin{lstlisting}[language=Matlab,caption={Configurazione principale e modello delle deadline.},label={lst:impl-configurazione}]
	numberOfResources = 8;
	averageT          = 180;
	maxT              = 8*3600;
	threshold         = Inf;
	showPlots         = true;
	
	deadlineModel.CellsPerSecond          = 5.5;
	deadlineModel.TravelMarginRatio       = 0.20;
	deadlineModel.HandlingSeconds         = 60;
	deadlineModel.MarginSecondsPerMission = 0;
	deadlineModel.DistanceMetric          = "euclidean";
	
	workFlow = importFromExcel( ...
	fullfile(projectRoot,'excel_files','1su11.xlsx'));
	
	[map,map_clear,costMap,signFollowerMap,occupancyMatrix] = ...
	computeMap( ...
	fullfile(projectRoot,'pathPlan','map1_1Filled.png'), ...
	fullfile(projectRoot,'pathPlan','map1_1.png'));
\end{lstlisting}

\section{Importazione e preparazione delle missioni}
\label{sec:impl-importazione}

\subsection{Lettura del foglio e normalizzazione delle colonne}

\texttt{import\allowbreak From\allowbreak Excel} legge il contenuto del foglio mediante
\texttt{xlsread}, interpreta la prima riga come intestazione e trasforma i
nomi delle colonne in maiuscolo eliminando gli spazi esterni. Sono richiesti
il codice missione e le quattro coordinate FROM e TO. La conversione dei
relativi valori presuppone una rappresentazione numerica coerente; non è
implementato un validatore generale di qualsiasi schema Excel.

Per l'identificativo dell'ordine viene cercata la prima colonna disponibile
tra \texttt{ID\_\allowbreak ORDER}, \texttt{OID\_\allowbreak ORDER} e \texttt{CD\_\allowbreak ORDER}, in questo
ordine. I valori della colonna scelta vengono convertiti in stringhe ed
esposti internamente come \texttt{ID\_\allowbreak ORDER}. La ricerca seleziona una
colonna per l'intero foglio: una cella vuota nella prima colonna disponibile
non viene riempita con il valore di una colonna alternativa.

In assenza di una colonna d'ordine riconosciuta viene emesso un avviso.
Gli identificativi mancanti vengono normalizzati a stringhe vuote, mentre
le missioni corrispondenti restano attività indipendenti. La priorità è
ricavata da \texttt{NM\_\allowbreak PRIORITY}; se la colonna o un valore sono mancanti,
si usa la priorità unitaria.

\subsection{Schema e identità delle missioni}

\texttt{define\allowbreak Missions} costruisce una tabella con i campi
\texttt{mission\allowbreak Code}, \texttt{priority}, \texttt{mission\allowbreak From},
\texttt{mission\allowbreak To}, \texttt{order\allowbreak Id} e \texttt{mission\allowbreak Uid}.
Gli estremi FROM e TO sono sottotabelle con le rispettive coordinate;
l'UID è un identificativo numerico positivo e univoco della singola
occorrenza. Se non è già presente viene generato dalla posizione della
riga, prima del filtraggio e dell'ordinamento.

Questa distinzione è necessaria perché il codice commerciale può ripetersi.
Due righe con lo stesso \texttt{mission\allowbreak Code} non vengono automaticamente
fuse in una sola attività. L'UID permette di aggiornare la riga corretta
nello stato delle deadline anche quando la lista viene riordinata.

\subsubsection{Costruzione della tabella delle missioni}

Il codice seguente mostra come l'identità stabile venga creata prima del
filtraggio. L'UID non dipende quindi dal successivo ordinamento per ordine
e priorità.

\begin{lstlisting}[language=Matlab,caption={Costruzione dei campi principali e dell'identificativo univoco.},label={lst:impl-define-missions}]
	missionFrom = workFlow(:, {'FROM_X','FROM_Y'});
	missionTo   = workFlow(:, {'TO_X','TO_Y'});
	
	priority = workFlow.NM_PRIORITY;
	priority(ismissing(priority)) = 1;
	
	missionCode = workFlow.CD_MISSION;
	orderId = string(workFlow.ID_ORDER);
	
	if ismember('missionUid',workFlow.Properties.VariableNames)
	missionUid = workFlow.missionUid;
	else
	missionUid = (1:height(workFlow))';
	end
	
	validateattributes(missionUid, ...
	{'numeric'},{'column','positive','integer','finite'});
	
	assert(numel(unique(missionUid)) == height(workFlow), ...
	'Mission UIDs must be unique.');
	
	missionList = table( ...
	missionCode,priority,missionFrom,missionTo,orderId,missionUid);
\end{lstlisting}

\subsection{Correzione e ammissibilità delle coordinate}

La preparazione applica dapprima l'euristica di correzione della scala:
le coordinate maggiori di 999 vengono divise per 1000 e arrotondate.
La regola deriva dalla rappresentazione dei dati disponibili e deve essere
verificata per nuovi fogli, perché non distingue un errore di importazione
da una coordinata realmente distante.

Successivamente ogni coordinata finita esterna al dominio viene riportata
entro i limiti della mappa. I valori non finiti restano invalidi. La funzione
restituisce due tabelle: \texttt{mission\allowbreak List}, contenente le attività
ammesse, e \texttt{discarded\allowbreak Missions}, contenente quelle escluse.
La dicitura stampata relativa alle missioni fuori mappa va letta alla luce
di questa correzione: una coordinata finita esterna può essere ammessa dopo
la proiezione, mentre una coordinata non finita non viene recuperata.

Le attivit? escluse non generano sottoreti e non vengono inserite nello
stato iniziale delle deadline del flusso corrente. Il loro conteggio
è quindi separato dal completamento degli ordini nel riepilogo finale.
Se la tabella delle missioni ammesse è vuota, lo script emette un avviso e
termina prima di avviare il motore.

\subsubsection{Validazione geometrica}

La separazione fra attività ammesse ed escluse è implementata mediante una
maschera logica costruita sulle quattro coordinate.

\begin{lstlisting}[language=Matlab,caption={Validazione delle coordinate in \texttt{defineMissions}.},label={lst:impl-validazione-coordinate}]
	xy = workFlow{:, {'FROM_X','FROM_Y','TO_X','TO_Y'}};
	
	finiteCoordinates = all(isfinite(xy),2);
	
	positiveCoordinates = ...
	xy(:,1)>=1 & xy(:,2)>=1 & ...
	xy(:,3)>=1 & xy(:,4)>=1;
	
	insideMap = ...
	xy(:,1)<=maxX & xy(:,2)<=maxY & ...
	xy(:,3)<=maxX & xy(:,4)<=maxY;
	
	valid = finiteCoordinates & ...
	positiveCoordinates & insideMap;
	
	discardedMissions = missionList(~valid,:);
	missionList       = missionList(valid,:);
\end{lstlisting}

\subsection{Raggruppamento stabile e priorità}

\texttt{group\allowbreak Missions\allowbreak By\allowbreak Order} assegna agli ordini il rango della loro prima
comparsa, crea gruppi distinti per le missioni senza ordine e ordina secondo
tre chiavi: gruppo, priorità decrescente e posizione originale.
Si conservano così sia l'ordine tra gli ordini sia la stabilità delle
missioni con priorità uguale.

La priorità non è un criterio globale per scegliere tra ordini diversi.
Determina la sequenza interna dell'ordine, che sarà imposta anche dai posti
di collegamento della rete. Il raggruppamento viene applicato prima della
suddivisione e nuovamente al batch corrente dopo eventuali reinserimenti.

\subsubsection{Ordinamento stabile nel codice}

La funzione assegna un gruppo distinto anche alle missioni prive di ordine.
La terza colonna fornita a \texttt{sortrows} conserva la posizione originaria
come criterio di spareggio.

\begin{lstlisting}[language=Matlab,caption={Raggruppamento stabile per ordine e priorità.},label={lst:impl-group-orders}]
	ids = strip(string(missions.orderId));
	[~,~,groups] = unique(ids,'stable');
	
	missing = ismissing(ids) | ids=="" | strcmpi(ids,"NaN");
	groups(missing) = max(groups) + (1:nnz(missing));
	
	[~,~,groups] = unique(groups,'stable');
	
	missions.priority(ismissing(missions.priority)) = 1;
	
	[~,indices] = sortrows( ...
	[groups,-missions.priority,(1:height(missions))']);
	
	missions = missions(indices,:);
\end{lstlisting}

\section{Assegnazione offline e batch}
\label{sec:impl-assegnazione}

\subsection{Compatibilità e proprietario dell'ordine}

\texttt{define\allowbreak Attitude} costruisce i limiti geometrici delle risorse.
Nella configurazione corrente questi limiti coincidono per tutte le risorse.
\texttt{assign\allowbreak Mission} controlla l'appartenenza delle coordinate FROM al
rettangolo operativo di ciascuna risorsa. Il controllo non riguarda
l'esistenza di un percorso e non verifica separatamente il TO come vincolo
di compatibilità: la fattibilità spaziale viene valutata dal pianificatore.

\texttt{assign\allowbreak Orders} considera tutte le missioni di ogni ordine e calcola
l'intersezione delle risorse compatibili. Se l'intersezione è vuota solleva
un errore, senza distribuire parzialmente l'ordine. Altrimenti sceglie la
risorsa con il minor numero di missioni già attribuite; a parità prevale
l'indice più piccolo.

La risorsa scelta viene scritta nel campo \texttt{assigned\allowbreak Resource} di
tutte le missioni dell'ordine. Una tabella separata,
\texttt{selected\allowbreak Order\allowbreak Assignments}, registra ordine, proprietario e
cardinalità ammessa. L'assegnazione avviene sull'intera lista prima di
costruire i batch e non dipende dalle durate A*, ancora non calcolate.

Quando \texttt{assign\allowbreak Mission} viene richiamata durante l'esecuzione,
riconosce il proprietario già presente, ne controlla la compatibilità e
restringe i colori ammessi a quel solo proprietario. Non effettua quindi
una nuova scelta libera tra risorse a ogni batch.

\subsubsection{Scelta del proprietario nel codice}

Il Listing~\ref{lst:impl-assign-orders} implementa direttamente
l'intersezione degli insiemi di risorse ammissibili e l'euristica
\emph{least-loaded}. Il carico usato in questa fase è il numero di missioni
già assegnate, non il tempo A*.

\begin{lstlisting}[language=Matlab,caption={Assegnazione univoca dell'ordine alla risorsa meno caricata.},label={lst:impl-assign-orders}]
	load = zeros(1,numberOfResources);
	
	for g = 1:max([0;groups])
	
	rows = find(groups == g);
	candidates = 1:numberOfResources;
	
	for k = rows'
	candidates = intersect(candidates,eligible{k});
	end
	
	if isempty(candidates)
	error('assignOrders:NoCommonResource', ...
	['Order %s has no resource eligible ' ...
	'for all its missions.'],ids(rows(1)));
	end
	
	[~,leastLoaded] = min(load(candidates));
	owner = candidates(leastLoaded);
	
	missions.assignedResource(rows) = owner;
	load(owner) = load(owner) + numel(rows);
	end
\end{lstlisting}

\subsection{Cardinalità nominale e integrità dell'ordine}

\texttt{batch\allowbreak Division} imposta la cardinalità nominale a $2R$, dove $R$
è il numero di risorse. Con otto risorse, il riferimento è dunque di 16
missioni per batch. Se il limite cade all'interno di un ordine, il batch
viene esteso fino all'ultima missione di quell'ordine. L'ultimo batch può
contenere meno missioni del riferimento.

La sequenza dei batch è conservata in un array di celle di tabelle,
\texttt{divided\allowbreak Missions}. Un ordine numeroso può produrre un batch maggiore
della dimensione nominale: il vincolo di integrità prevale sulla cardinalità.
La suddivisione è determinata prima della simulazione e non deriva da un
confronto tra configurazioni candidate.

\subsubsection{Integrazione nel flusso principale}

L'assegnazione viene completata prima di invocare il batching. Di conseguenza
ogni tabella contenuta in \texttt{dividedMissions} eredita già il campo
\texttt{assignedResource}.

\begin{lstlisting}[language=Matlab,caption={Assegnazione degli ordini e successiva divisione in batch.},label={lst:impl-main-assignment-batch}]
	attitudeTable = defineAttitude( ...
	numberOfResources,size(costMap));
	
	[runMissionList,selectedOrderAssignments] = ...
	assignOrders( ...
	missionList,numberOfResources,attitudeTable);
	
	deadlineState = initializeDeadlineState(runMissionList);
	
	[dividedMissions,~] = batchDivision( ...
	runMissionList,numberOfResources, ...
	averageT,maxT,false);
\end{lstlisting}

\section{Costruzione della rete colorata}
\label{sec:impl-rete}

\subsection{Matrici di incidenza e indicizzazione}

\texttt{generate\allowbreak Matrices} associa a ogni missione tre transizioni e due
posti interni. Le matrici di pre-incidenza e post-incidenza hanno le
transizioni sulle righe e i posti sulle colonne; questa convenzione va
considerata nell'interpretazione degli indici.

Per la missione in riga $i$, le transizioni sono
$t_{3i-2}$, $t_{3i-1}$ e $t_{3i}$, corrispondenti all'avvio, alla fase
con durata e al passaggio finale. Oltre ai posti interni viene creato il
posto comune $p_1$ delle risorse disponibili. Per un batch con $M$ missioni
e $H$ collegamenti interni agli ordini, le dimensioni sono
\begin{equation}
	N_T=3M,\qquad N_P=1+2M+H.
	\label{eq:impl-dimensioni-rete}
\end{equation}
Un ordine con $q$ missioni ammesse introduce $q-1$ collegamenti; una
missione senza ordine resta una catena autonoma senza tali collegamenti.

Per collegare due missioni consecutive, la funzione elimina il ritorno
della precedente a $p_1$ e il prelievo della successiva da $p_1$, sostituendoli
con un posto di passaggio. Il token della risorsa rimane nella catena fino
all'ultima missione. La sequenza viene ricostruita per priorità anche se
le righe ricevute non sono già ordinate.

\subsection{Archi, colori e marcatura}

\texttt{generate\allowbreak Arcs} traduce gli elementi non nulli delle matrici in una
tabella con origine, destinazione e matrice dei colori ammessi.
I nomi dei nodi distinguono posti e transizioni; il collegamento tra
transizione e missione si ricava dal gruppo di tre transizioni cui essa
appartiene.

Ogni risorsa è rappresentata da un vettore di base di lunghezza $R$,
con un solo elemento unitario. Dopo l'assegnazione offline, gli archi
della missione ammettono esclusivamente il colore del proprietario.
Questo vale anche per gli archi dei posti di passaggio.

La marcatura è una matrice $N_P\times R$. Il programma usa il numero di
posti effettivamente generato, compresi quelli di collegamento; inizializza
la prima riga a uno per tutte le risorse e le altre a zero. Non vengono
aggiunti token di controllo separati per le precedenze. L'invariante
operativo è un token per colore, distribuito tra i posti della rete:
\begin{equation}
	\sum_{p=1}^{N_P}m_{p,r}=1,\qquad r=1,\ldots,R.
	\label{eq:impl-token}
\end{equation}

\subsubsection{Generazione di rete, archi e marcatura}

Nel ciclo principale le strutture vengono rigenerate per il batch corrente.
La dimensione della marcatura è ricavata direttamente da
\texttt{size(pre\_m,2)}, così da includere anche gli eventuali posti di
handoff aggiunti dalla rete.

\begin{lstlisting}[language=Matlab,caption={Costruzione della CTPN del batch e marcatura iniziale.},label={lst:impl-build-pn}]
	[~,colors] = assignMission( ...
	currentBatch,numberOfResources,attitudeTable);
	
	[pre_m,post_m] = generateMatrices(currentBatch);
	arc = generateArcs(pre_m,post_m,colors);
	
	if ~iscell(arc.color_mat)
	arc.color_mat = mat2cell( ...
	arc.color_mat, ...
	ones(1,size(arc,1)), ...
	numberOfResources);
	end
	
	m = zeros(size(pre_m,2),numberOfResources);
	m(1,:) = ones(1,numberOfResources);
\end{lstlisting}

\section{Motore di esecuzione e stato conservato}
\label{sec:impl-stato}

\subsection{Abilitazione e scatto delle transizioni}

\texttt{colored\allowbreak Petri} mantiene l'elenco dei posti da elaborare, gli archi
residui e la cronologia degli scatti. Quando un posto alimenta più
transizioni, \texttt{r\_\allowbreak conflict\_\allowbreak colored} individua quelle compatibili
con la marcatura e con i colori ammessi. Le disponibilità vengono aggiornate
anche durante la selezione, evitando di impegnare due volte lo stesso token.

La routine conserva un criterio storico basato sul lavoro accumulato per
scegliere tra colori disponibili. Nel flusso corrente ogni missione ammette
un solo proprietario: tale criterio non può modificarlo. Per un'uscita
singola interviene \texttt{check\_\allowbreak color}, che verifica la disponibilità del
colore richiesto.

\texttt{next\_\allowbreak transition\_\allowbreak c} determina il tempo associato al token in
ingresso mediante \texttt{find\_\allowbreak position}. Quest'ultima ricerca nella
cronologia l'occorrenza più recente del nodo con il colore pertinente;
quando necessario riconosce la marcatura iniziale con tutte le risorse.
La transizione centrale della missione richiede la pianificazione del
percorso e riceve la durata calcolata. Le transizioni di avvio e passaggio
non aggiungono un ulteriore tempo operativo.

\texttt{next\_\allowbreak place\_\allowbreak c} trasferisce il colore ai posti in uscita e propaga
il tempo della transizione. Gli archi attraversati vengono rimossi dalla
tabella residua. La cronologia restituita, \texttt{schedule\_\allowbreak sort}, viene
ordinata per tempo e filtrata per conservare le transizioni.

L'ordine di elaborazione delle strutture nel programma non va confuso con
una scansione globale dell'orologio a intervalli fissi. Il motore costruisce
i tempi delle attività propagandoli lungo i token e produce successivamente
la cronologia ordinata.

\subsection{Avvio, risultato e posizione della risorsa}

Per una missione con pianificazione riuscita, il tempo del token in ingresso
alla fase centrale è l'avvio $s_m$. Il completamento è
\begin{equation}
	c_m=s_m+\tau_m.
	\label{eq:impl-completamento}
\end{equation}
Durante questa elaborazione vengono già registrati il risultato completato,
i tempi previsti dal percorso e il nuovo stato della risorsa. Non si attende
il trascorrere di un tempo reale sul computer e non si conserva una missione
in esecuzione con un residuo da riprendere alla finestra successiva.

\texttt{update\allowbreak Batch} aggiunge la durata al campo \texttt{Time} della
risorsa, appende codice, ordine e priorità allo storico e pone la posizione
sull'ultimo punto del percorso. La tabella \texttt{resource\allowbreak Pos} contiene
una matrice numerica $R\times2$ nel campo \texttt{position}, inizializzata
con valori non definiti. Il campo \texttt{resource} contiene i vettori
unitari che permettono di riconoscere la risorsa dal colore.

I tempi del programma di esecuzione rimangono simulati. Il loro aggiornamento
immediato durante il calcolo non rappresenta un aggiornamento anticipato
di un dispositivo fisico, che non fa parte del sistema implementato.

\subsection{Passaggio al batch successivo}

Al termine di un batch, \texttt{main.m} conserva il riepilogo cumulato
in \texttt{scheduled\allowbreak Mission}, le posizioni e \texttt{deadline\allowbreak State}.
Le copie dei riepiloghi e delle cronologie sono memorizzate in
\texttt{batch\allowbreak History} e \texttt{schedule\allowbreak History}. La matrice
\texttt{batch\allowbreak Time} contiene i tempi cumulati in ore, con una colonna per
batch; non rappresenta il solo incremento di lavoro dell'ultimo batch.

Nel batch successivo rete, archi e marcatura sono nuovi. Il motore associa
al token di ciascun proprietario il suo precedente tempo cumulato. Non
impone a tutte le risorse di ripartire dal massimo dei tempi raggiunti dalle
altre: il passaggio tra batch non introduce una barriera di sincronizzazione.

Il controllo della soglia resta presente, ma con valore infinito non blocca
le risorse nel normale scenario. I rami storici che bloccano una risorsa per
bilanciare il lavoro sono distinti dal ramo usato quando esiste già un
proprietario assegnato.

\subsubsection{Ciclo di simulazione dei batch}

Il Listing~\ref{lst:impl-main-loop} mostra il punto in cui il motore riceve
la rete appena costruita e restituisce contemporaneamente cronologia,
lavoro cumulato, posizioni, missioni rinviate, fallimenti e stato delle
deadline.

\begin{lstlisting}[language=Matlab,caption={Invocazione del motore e persistenza dello stato fra batch.},label={lst:impl-main-loop}]
	[schedule_sort,batch,resourcePos, ...
	notDoneMissions,failedMissionCodes,deadlineState] = ...
	coloredPetri( ...
	pre_m,post_m,m,numberOfResources,arc, ...
	scheduledMission,resourcePos,k, ...
	planner,occupancyMatrix,currentBatch,threshold, ...
	deadlineState,deadlineModel,runMissionList);
	
	scheduledMission = batch;
	
	scheduleHistory{k} = schedule_sort;
	batchHistory{k}    = batch;
	batchTime(:,end+1) = batch.Time/3600;
	
	if ~isempty(failedMissionCodes)
	failedRows = currentBatch( ...
	ismember(currentBatch.missionCode,failedMissionCodes),:);
	
	failedPathMissions = ...
	[failedPathMissions; failedRows];
	end
\end{lstlisting}

\section{Mappa e pianificazione dei percorsi}
\label{sec:impl-percorsi}

\subsection{Costruzione della griglia}

\texttt{compute\allowbreak Map} legge l'immagine degli ostacoli e quella di riferimento.
Per immagini a colori viene usato il canale verde. I pixel dell'immagine
degli ostacoli diversi da 255 ricevono costo 100, mentre quelli bianchi
rimangono a costo zero. Questi valori alimentano la mappa di occupazione
usata dal pianificatore.

La costruzione di \texttt{binary\allowbreak Occupancy\allowbreak Map} usa esplicitamente dimensioni
50.2 per 95.1 e risoluzione 10, coerenti con la griglia di riferimento di
951 righe e 502 colonne. Sostituire l'immagine con una mappa di dimensioni
diverse richiede quindi di verificare anche questa configurazione.
Il nome \texttt{cost\allowbreak Map} non indica, nel flusso corrente, un modello di
costo della congestione o di interferenza tra mezzi.

\subsection{Trasformazione degli estremi}

\texttt{scale\allowbreak Points} trasforma le coordinate operative in indici di
matrice secondo
\begin{equation}
	(X,Y)\longmapsto
	\bigl(\operatorname{round}(X),
	\operatorname{round}(N_Y+1-Y)\bigr).
	\label{eq:impl-coordinate}
\end{equation}
Il termine unitario mantiene la validità degli indici MATLAB ai bordi:
$Y=1$ corrisponde all'ultima colonna e $Y=N_Y$ alla prima.
Le coordinate sono quindi arrotondate e l'asse Y invertito.

\texttt{path\allowbreak Planning} legge l'occupazione direttamente dalla mappa del
planner e verifica la coerenza dimensionale con la matrice ricevuta.
Imposta la connettività della ricerca di raggiungibilità a quattro oppure
otto vicini in base alla proprietà \texttt{Diagonal\allowbreak Search} del planner.
Lo script principale non modifica esplicitamente tale proprietà.

\subsection{Continuità e correzione degli estremi}

Ogni chiamata pianifica due tratte: posizione corrente--FROM e FROM--TO.
Se la risorsa non ha una posizione precedente, viene inizializzata sul
FROM richiesto. Non si introduce un trasferimento da un deposito comune.
Dopo il controllo degli indici, una posizione iniziale occupata viene
sostituita dalla cella libera più vicina.

Un FROM occupato viene cercato nella componente raggiungibile dalla
posizione iniziale. Un FROM libero ma isolato non viene invece spostato:
la prima chiamata ad A* può fallire. Se l'avvicinamento riesce, il suo ultimo
punto diventa il FROM effettivo della seconda tratta, evitando correzioni
indipendenti che interromperebbero la continuità.

Il TO viene sempre confrontato con la componente raggiungibile dal FROM
effettivo. Se è libero e appartiene a tale componente rimane invariato;
altrimenti viene sostituito dalla cella della componente più vicina al TO
richiesto. La regola usa la distanza euclidea al quadrato e risolve le parità
per riga e poi per colonna.

Una volta calcolata la seconda tratta, le due sequenze vengono concatenate.
Il percorso è valido rispetto agli ostacoli statici considerati dal planner;
non viene confrontato con le occupazioni temporali degli altri mezzi.
La correzione di un estremo garantisce soltanto una scelta geometrica e non
la corrispondenza con la postazione logistica originaria.

\subsection{Calcolo della durata simulata}

Indicando con $K_m$ il numero di righe della sequenza concatenata, il tempo
restituito da \texttt{path\allowbreak Planning} è
\begin{equation}
	\tau_m=
	\operatorname{round}\!\left(\frac{K_m}{5.5}\,1.20\right)+60
	\quad [\mathrm{s}].
	\label{eq:impl-durata}
\end{equation}
Il conteggio comprende gli estremi delle due tratte e il punto FROM presente
in entrambe. La formula usa il numero di punti, non la somma delle lunghezze
degli archi: i movimenti diagonali non ricevono un peso geometrico separato
nel calcolo temporale.

I coefficienti della durata sono scritti nella funzione di pianificazione.
La previsione delle deadline riceve invece i propri coefficienti dalla
struttura di configurazione. Nella versione corrente i valori coincidono,
ma cambiare la configurazione delle deadline non modifica automaticamente
la legge temporale del percorso. Questa separazione deve essere considerata
quando si calibrano i parametri.

\subsubsection{Implementazione del percorso e del tempo}

Il nucleo della funzione \texttt{pathPlanning} rende esplicite le due tratte
della missione e la successiva conversione del percorso in secondi. Il
Listing~\ref{lst:impl-path-planning} riporta l'estratto corrispondente.

\begin{lstlisting}[language=Matlab,caption={Pianificazione delle due tratte e calcolo del tempo simulato.},label={lst:impl-path-planning}]
	path = plan(planner,startPoint,fromPoint);
	
	if isempty(path)
	fprintf(['Discarded mission %d: ' ...
	'no feasible approach path.\n'], ...
	mission.missionCode);
	return
	end
	
	% L'ultimo punto raggiunto diventa il FROM effettivo.
	fromPoint = path(end,:);
	
	secondPath = plan(planner,fromPoint,toPoint);
	
	if isempty(secondPath)
	fprintf(['Discarded mission %d: ' ...
	'no feasible FROM-to-TO path.\n'], ...
	mission.missionCode);
	return
	end
	
	path = [path;secondPath];
	
	timePath = size(path,1)/5.5;
	time = round(timePath + timePath*0.2) + 60;
	
	missionTable = table( ...
	mission.missionCode, ...
	mission.priority, ...
	time,{path}, ...
	'VariableNames', ...
	{'MissionCode','Priority','Time','Path'});
\end{lstlisting}

\section{Riutilizzo delle componenti raggiungibili}
\label{sec:impl-cache}

\texttt{find\allowbreak Nearest\allowbreak Zero} gestisce sia la ricerca della cella libera più
vicina sia quella vincolata alla componente di un'origine. Nel secondo
caso conserva dati persistenti: maschera delle celle libere, connettività,
etichette delle componenti e coordinate dei relativi punti.

Una componente ancora sconosciuta viene esplorata con una visita in ampiezza.
Le chiamate successive dalla stessa componente riutilizzano il risultato.
Se il punto richiesto risulta già etichettato come appartenente alla
componente corretta, viene restituito direttamente. Negli altri casi si
calcolano le distanze dai punti della componente e si applica il criterio
deterministico di parità.

La cache viene invalidata quando cambia la maschera delle celle libere o
la connettività, anche se le dimensioni della matrice restano uguali.
Riguarda la mappa più recente e non conserva i percorsi A* delle singole
missioni. Una chiamata verifica comunque la coerenza della maschera; il
riutilizzo non rende ogni ricerca un'operazione a costo costante.

Il vantaggio consiste nell'evitare visite ripetute della medesima componente.
Non modifica gli estremi scelti, le precedenze o la durata ottenuta da un
identico percorso. Una riduzione del tempo di calcolo non va pertanto
interpretata come una riduzione del lavoro fisico simulato.

\section{Implementazione delle deadline}
\label{sec:impl-deadline}

\subsection{Stato iniziale di missioni e ordini}

\texttt{initialize\allowbreak Deadline\allowbreak State} riceve la lista ammessa e assegnata.
Costruisce due tabelle, \texttt{Missions} e \texttt{Orders}, con gli stati
iniziali in attesa e i valori temporali non ancora definiti.
Per gli ordini identificati usa una chiave derivata dall'ID; per una missione
senza ordine usa una chiave basata sull'UID. Le missioni senza ID non vengono
quindi riunite in un unico ordine.

La tabella delle missioni conserva UID, codice, chiave e nome dell'ordine,
risorsa, posizione nella sequenza, distanza e durata stimate, deadline,
avvio, completamento, durata simulata, slack, ritardo positivo e stato.
Quella degli ordini contiene proprietario, cardinalità ammessa, previsione
complessiva e indicatori temporali aggregati. L'inizializzazione verifica
che ogni ordine abbia un unico proprietario.

\subsection{Congelamento al primo avvio riuscito}

La chiamata a \texttt{freeze\allowbreak Order\allowbreak Deadline} avviene nella transizione centrale,
solo dopo che \texttt{path\allowbreak Planning} ha restituito un percorso valido e prima
dell'aggiornamento della posizione. La stima riceve quindi la posizione
precedente della risorsa e il tempo del token all'avvio.

La funzione individua l'ordine tramite l'UID della missione e seleziona dalla
lista globale tutte le sue missioni ammesse. Se il tempo iniziale dell'ordine
è già definito, termina senza alterare la previsione. Diversamente richiama
\texttt{estimate\allowbreak Order\allowbreak Deadlines\allowbreak Euclidean}, registra tutte le scadenze
cumulative e pone l'ordine nello stato in esecuzione.

Il riferimento viene congelato una sola volta. Non viene traslato in avanti
per assorbire ritardi successivi. Se le prime missioni falliscono prima di
un avvio riuscito, la previsione comprende comunque l'intera lista ammessa
dell'ordine, incluse quelle fallite; se tutte falliscono, avvio e deadline
restano indefiniti.

\subsubsection{Integrazione dello stato delle deadline}

Nel flusso attivo lo stato viene creato dopo l'assegnazione degli ordini e
viene poi passato per valore restituito a ogni chiamata del motore. Il
Listing~\ref{lst:impl-deadline-flow} evidenzia questo collegamento senza
introdurre un secondo orologio di simulazione.

\begin{lstlisting}[language=Matlab,caption={Persistenza di \texttt{deadlineState} nel flusso principale.},label={lst:impl-deadline-flow}]
	deadlineState = initializeDeadlineState(runMissionList);
	
	% ... costruzione del batch e della rete ...
	
	[schedule_sort,batch,resourcePos, ...
	notDoneMissions,failedMissionCodes,deadlineState] = ...
	coloredPetri( ...
	pre_m,post_m,m,numberOfResources,arc, ...
	scheduledMission,resourcePos,k, ...
	planner,occupancyMatrix,currentBatch,threshold, ...
	deadlineState,deadlineModel,runMissionList);
\end{lstlisting}

\subsection{Previsione geometrica per missione}

La funzione di stima ordina le missioni per priorità e applica la stessa
trasformazione degli indici di griglia. Non richiama A* né la ricerca della
cella raggiungibile. Per la missione $i$, con posizione prevista $Q_i$ e
estremi richiesti $F_i,T_i$, calcola
\begin{equation}
	\widehat\ell_i=\lVert Q_i-F_i\rVert_2+
	\lVert F_i-T_i\rVert_2.
	\label{eq:impl-distanza-deadline}
\end{equation}
Per la prima missione $Q_1$ coincide con la posizione ricevuta; se essa non
è definita, coincide con $F_1$. Per le successive si usa $Q_i=T_{i-1}$,
senza applicare le eventuali correzioni che il planner adotterà realmente.

La durata stimata è
\begin{equation}
	\widehat\tau_i=
	\operatorname{round}\!\left(
	\frac{\widehat\ell_i}{v_c}(1+\alpha)\right)
	+\tau_f+\delta_m,
	\label{eq:impl-stima-durata}
\end{equation}
con i parametri configurati nello script principale. Il margine
$\delta_m$ si applica a ogni missione ed è attualmente nullo.
L'arrotondamento avviene prima della somma, separatamente per missione.

Dato l'avvio $S_o$, le scadenze vengono accumulate secondo
\begin{equation}
	d_{o,i}=S_o+\sum_{j=1}^{i}\widehat\tau_j,
	\qquad
	d_o=S_o+\sum_{j=1}^{n_o^{\mathrm{amm}}}\widehat\tau_j.
	\label{eq:impl-deadline-cumulativa}
\end{equation}
L'ultima scadenza di missione coincide con quella dell'ordine. Il campo della
metrica e le funzioni di registrazione hanno finalità osservative: non
modificano la marcatura o i colori ammessi.

\subsection{Risultati, slack e classificazione}

\texttt{record\allowbreak Mission\allowbreak Deadline\allowbreak Result} identifica la missione tramite UID.
In caso di completamento memorizza avvio, fine e durata, quindi calcola
lo slack come deadline meno completamento e il ritardo positivo come
massimo tra zero e completamento meno deadline. Un esito fallito aggiorna
lo stato senza introdurre tempi di lavoro fittizi.

L'ordine è completato quando tutte le sue missioni ammesse sono completate;
la sua fine è il massimo dei tempi di completamento. Quando tutte le
missioni sono terminate ma almeno una è fallita, risulta incompleto.
Se è stato avviato e restano missioni non terminate, rimane in esecuzione.
Queste classificazioni sono riferite alla lista ammessa, perché le esclusioni
preliminari non sono state inserite nello stato.

Il ritardo viene calcolato sul risultato completato, non aggiornato
continuamente rispetto a un orologio globale. Un ordine incompleto non
viene contato tra quelli completati entro la deadline. Non è inoltre
previsto nel flusso attivo un limite manuale per ordine che sostituisca
la previsione euclidea.

\subsubsection{Estrazione degli indicatori finali}

La classificazione finale utilizza direttamente la tabella
\texttt{deadlineState.Orders}. Un ordine viene contato come in ritardo solo
se risulta completato e presenta \texttt{LateBy} strettamente positivo.

\begin{lstlisting}[language=Matlab,caption={Riepilogo finale degli ordini rispetto alle deadline.},label={lst:impl-deadline-summary}]
	missionDeadlineSummary = sortrows( ...
	deadlineState.Missions, ...
	{'Resource','OrderId','Sequence'});
	
	orderDeadlineSummary = sortrows( ...
	deadlineState.Orders, ...
	{'Resource','StartTime'});
	
	completedOrders = ...
	orderDeadlineSummary.Status == "completed";
	
	lateOrders = ...
	completedOrders & orderDeadlineSummary.LateBy > 0;
	
	onTimeOrders = ...
	completedOrders & orderDeadlineSummary.LateBy == 0;
	
	fprintf(['Ordini completati: %d | ' ...
	'entro deadline: %d | in ritardo: %d\n'], ...
	nnz(completedOrders), ...
	nnz(onTimeOrders), ...
	nnz(lateOrders));
\end{lstlisting}


\section{Risultati preliminari della simulazione}
\label{sec:risultati-preliminari-simulazione}

Per verificare il comportamento complessivo dell'architettura implementata è
stata eseguita una simulazione completa utilizzando la configurazione descritta
nelle sezioni precedenti, con otto risorse e con il motore
\texttt{coloredPetri} come nucleo dell'esecuzione. Lo scopo di questa prova non
è ancora effettuare un confronto sistematico tra differenti politiche di
scheduling, ma verificare che la catena completa
\[
\text{importazione}
\rightarrow
\text{assegnazione}
\rightarrow
\text{batching}
\rightarrow
\text{CTPN}
\rightarrow
\text{path planning}
\rightarrow
\text{deadline}
\]
produca risultati coerenti e osservabili.

\subsection{Indicatori globali}

La simulazione MATLAB è stata completata in $197.807$ s. Tutti i $310$ ordini
presenti nella prova sono stati portati a completamento. Non sono state
registrate missioni con percorso non fattibile e, al termine dell'ultimo batch,
non risultano missioni residue da rinviare a una successiva esecuzione.

I principali indicatori sono riassunti nella
Tabella~\ref{tab:simulation-global-results}.

\begin{table}[htbp]
	\centering
	\caption{Indicatori globali della simulazione di riferimento.}
	\label{tab:simulation-global-results}
	\begin{tabular}{lr}
		\toprule
		\textbf{Indicatore} & \textbf{Valore} \\
		\midrule
		Numero di risorse & 8 \\
		Ordini completati & 310 \\
		Tempo di calcolo MATLAB & 197.807 s \\
		Makespan simulato & 8.53 h \\
		Missioni con path non fattibile & 0 \\
		Missioni residue dopo l'ultimo batch & 0 \\
		Ordini entro deadline & 309 \\
		Ordini in ritardo & 1 \\
		Percentuale ordini entro deadline & 99.68\% \\
		Massimo ritardo osservato & 2 s \\
		\bottomrule
	\end{tabular}
\end{table}

Il valore del makespan deve essere interpretato rispetto alla formulazione
adottata nel simulatore. Il riferimento di otto ore non costituisce infatti
una condizione di arresto: le attività già assegnate possono proseguire oltre
tale valore. Nella prova considerata il massimo carico cumulato è pari a
$8.5303$ h ed è associato alla risorsa 4.

\subsection{Distribuzione del carico tra le risorse}

La Tabella~\ref{tab:simulation-resource-workload} riporta il tempo di lavoro
cumulato prodotto dal motore per ciascuna risorsa.

\begin{table}[htbp]
	\centering
	\caption{Tempo di lavoro cumulato per risorsa.}
	\label{tab:simulation-resource-workload}
	\begin{tabular}{cc}
		\toprule
		\textbf{Risorsa} & \textbf{Tempo di lavoro [h]} \\
		\midrule
		1 & 7.1392 \\
		2 & 6.3403 \\
		3 & 6.9453 \\
		4 & 8.5303 \\
		5 & 5.5744 \\
		6 & 7.2911 \\
		7 & 7.9036 \\
		8 & 7.0908 \\
		\bottomrule
	\end{tabular}
\end{table}

Il carico medio è pari a circa $7.10$ h per risorsa. La distribuzione non è
perfettamente uniforme: la risorsa 4 raggiunge il valore massimo, mentre la
risorsa 5 presenta il valore minimo, pari a $5.5744$ h. Questo risultato è
coerente con il fatto che l'assegnazione iniziale bilancia il numero di
missioni e rispetta i vincoli di compatibilità, ma non risolve un problema di
ottimizzazione globale dei tempi effettivi ottenuti successivamente mediante
A*.

La differenza osservata tra i carichi costituisce quindi un indicatore utile
per le successive estensioni dello scheduler: una futura politica dinamica può
utilizzare il lavoro già accumulato, insieme alle deadline, per modificare
l'assegnazione degli ordini residui.

\subsection{Verifica delle deadline}

La stima geometrica delle deadline ha prodotto $309$ ordini completati entro
il riferimento temporale e un solo ordine in ritardo. La percentuale di ordini
completati entro deadline è quindi
\[
\frac{309}{310}\cdot100 \simeq 99.68\%.
\]

La Tabella~\ref{tab:simulation-deadline-examples} riporta tre casi
rappresentativi, scelti per mostrare le tre possibili condizioni dello slack:
positivo, nullo e negativo.

\begin{table}[htbp]
	\centering
	\caption{Esempi di confronto tra completamento effettivo e deadline.}
	\label{tab:simulation-deadline-examples}
	\resizebox{\textwidth}{!}{%
		\begin{tabular}{lrrrrrrrl}
			\toprule
			\textbf{OrderId} &
			\textbf{Ris.} &
			\textbf{Missioni} &
			\textbf{Start [s]} &
			\textbf{$\widehat{D}_o$ [s]} &
			\textbf{Deadline [s]} &
			\textbf{Finish [s]} &
			\textbf{Slack [s]} &
			\textbf{LateBy [s]} \\
			\midrule
			2006259 & 1 & 1 & 0     & 231 & 231   & 208   & 23 & 0 \\
			2015531 & 7 & 1 & 10525 & 75  & 10600 & 10600 & 0  & 0 \\
			2017195 & 3 & 7 & 1383  & 552 & 1935  & 1937  & -2 & 2 \\
			\bottomrule
	\end{tabular}}
\end{table}

Il primo caso mostra un ordine completato con margine positivo. Il secondo
raggiunge esattamente la deadline e viene pertanto classificato come
completato entro il limite. Il terzo rappresenta l'unico ordine tardivo della
simulazione: l'ordine \texttt{2017195}, assegnato alla risorsa 3, termina a
$t=1937$ s rispetto a una deadline pari a $1935$ s. Di conseguenza,
\[
Slack=-2~\mathrm{s},
\qquad
LateBy=2~\mathrm{s}.
\]
Questo caso costituisce una verifica diretta della logica implementata: uno
slack negativo viene trasformato in un ritardo positivo della stessa entità,
senza modificare retroattivamente la deadline.

\begin{figure}[H]
	\centering
	\includegraphics[width=1\linewidth]{imgs/esempi/resourceLoad}
	\caption{Distribuzione del carico lavorativo complessivo su tutte le risorse.}
	\label{fig:resourceload}
\end{figure}

\begin{figure}[H]
	\centering
	\includegraphics[width=1\linewidth]{imgs/esempi/priorityOrder}
	\caption{Esempio di scheduling dei singoli ordini: è possibile vedere le missioni ordinate per priorità, dalla più alta a quella più bassa.}
	\label{fig:priorityorder}
\end{figure}


\begin{figure}[H]
	\centering
	\includegraphics[width=1\linewidth]{imgs/esempi/orderTimeline}
	\caption{Distribuzione degli ordini per risorsa compreso della visualizzazione delle singole risorse.}
	\label{fig:ordertimeline}
\end{figure}


\begin{figure}[H]
	\centering
	\includegraphics[width=1\linewidth]{imgs/esempi/orderTimeline2}
	\caption{Esempio di visualizzazione delle deadline dei singoli ordini compreso di esecuzione e scadenze.}
	\label{fig:ordertimeline2}
\end{figure}



\section{Dati sperimentali}

\subsection{Stima euclidea per il bilanciamento dei batch}

Per bilanciare il carico delle risorse, il controllore utilizza una
stima della durata delle missioni basata sulla distanza euclidea,
in sostituzione del calcolo preventivo dei percorsi tramite A*.
La stima considera lo spostamento dalla posizione prevista della
risorsa al punto di prelievo e da questo alla destinazione,
includendo un margine del 20\% sul tempo di percorrenza e
60 secondi di movimentazione per missione.
A* rimane utilizzato nella fase di esecuzione per determinare
il percorso effettivamente percorribile.

Nel confronto effettuato sullo stesso insieme di 1\,113 missioni,
il tempo di calcolo MATLAB è passato da circa 21,4 minuti,
con la valutazione preventiva tramite A*, a 3,5 minuti con
la stima euclidea, con una riduzione di circa l'84\%.
Il tempo simulato complessivo è passato da 9,36 a 9,23 ore,
mentre gli ordini in ritardo sono aumentati da 39 a 43.

In questa prova, la stima euclidea ha quindi consentito un
notevole risparmio computazionale, mantenendo un tempo
complessivo di esecuzione comparabile, a fronte di un lieve
aumento degli ordini in ritardo. Pur non considerando
esplicitamente gli ostacoli, risulta utile per decidere
quante missioni ammettere in ciascun batch, evitando di
calcolare i percorsi sia durante l'assegnazione sia durante
l'esecuzione.